\section{The complete early CUE collision process}
\label{sec:analytic-cue-collision-process}

The first-event result NN-A024 does not by itself control the
other zeros after its first pair becomes nonreal. The following
uniform local argument supplies that missing exhaustion step.
The heat normalization is unchanged:
$Z_{N,-a}=e^{a\partial_z^2}Z_N$ and
$Z_N(z)=\prod_{j=1}^N2\sin((z-\theta_j)/2)$.

\begin{lemma}[Uniform one-root and two-root tracking]
\label{lem:analytic-uniform-clusters}
Let $f$ be real entire of order at most one with all zeros real.
Remove either a simple zero $b$, or two adjacent simple zeros
$b\pm\delta/2$, and call the remaining entire factor $G$.
Put $L=G'(b)/G(b)$ and $S=\sum_{G(\xi)=0}(b-\xi)^{-2}$.
Possible linear exponential factors are included in $L$.
There are absolute positive constants $c,C$ with the following
properties for $T>0$.

If one zero was removed and $TS<c$, then for every $a\in[0,T]$
exactly one zero of $e^{a\partial_z^2}f$ lies in the moving disc
\[
                  |z-(b-2aL)|<2\sqrt T.
\]
It is real and simple throughout the interval.

If two zeros were removed, put $R^2=\delta^2+T$. If
$R^2S<c$, exactly two zeros lie in
$|z-(b-2aL)|<2R$ for every $a\in[0,T]$. They have at most
one collision there. The collision is double and crosses from
two real zeros to a nonreal conjugate pair; the pair cannot
return to the real axis during $[0,T]$. Its collision time,
when in this interval, is the time supplied by NN-A023:
\[
             a_c=\frac{\delta^2}{8}
                \bigl(1+\delta^2S/4+O((\delta^2S)^2)\bigr).
\]
The constants are independent of $L$, and of the ratio
$\delta^2/T$.
\end{lemma}

\begin{proof}
Use the canonical remainder
\[
 \mathcal R(h)=G(b+h)/(G(b)e^{Lh}),\qquad
 |\mathcal R(h)|\le e^{S|h|^2/2},\qquad
 |\mathcal R(h)-1|\le 8S|h|^2e^{2S|h|^2},
\]
proved in NN-A022. Conjugating the Gaussian heat operator by
$e^{L(z-b)}$ removes the common drift exactly. In the one-root
case, with $z=b-2aL+\sqrt T\,u$ and $\tau=a/T$, division by
$\sqrt T G(b)e^{\sqrt T uL-aL^2}$ gives
\[
 I(\tau,u)=\mathbb E\bigl[Y\mathcal R(\sqrt T Y)\bigr],
                    \qquad Y=u+\sqrt{2\tau}\,V,
                    \qquad V\sim N(0,1).
\]
On any fixed bounded complex $u,\tau$ set this is
$I=u+O(TS)$, with every derivative on smaller compact sets
bounded in the same way. To justify this including complex
$\tau$, the Gaussian expectation is even in its square-root
parameter. Its absolute integrand is dominated by
$C(1+|V|)^C e^{C TS|V|^2}$; for small $TS$ the usual Gaussian
density dominates it uniformly. The expression therefore
extends holomorphically through $\tau=0$, and the stated
bounds and their derivatives follow by Cauchy's formula.
Rouch\'e on $|u|=2$ gives exactly one zero. The derivative
$I_u=1+O(TS)$ there is nonzero at that zero. Real conjugation
and uniqueness make the zero real for real $\tau$.

For the two-root case scale by $R$ instead, put
$d=\delta/R\in(0,1]$, $\tau=a/R^2\in[0,1]$, and divide by
$R^2G(b)e^{RuL-aL^2}$. The resulting exact family is
\[
 J(\tau,u)=\mathbb E\bigl[(Y^2-d^2/4)\mathcal R(RY)\bigr]
            =u^2+2\tau-d^2/4+O(R^2S).
\]
The same Gaussian argument is uniform in $0<d\le1$ and
controls all derivatives on fixed compact sets. On $|u|=2$
for real $0\le\tau\le1$, the leading polynomial has modulus
at least $4-2-1/4=7/4$. Thus exactly two zeros remain inside.
Its derivative has exactly one zero there, near zero, and
$J_{uu}=2+O(R^2S)>0$ on the real portion of this disc.
The critical point $u_*(\tau)$ is real analytic. Define
$D(\tau)=J(\tau,u_*(\tau))$. Because $J_u=0$ at $u_*$,
\[
                         D'(\tau)=J_\tau(\tau,u_*)
                                      =2+O(R^2S)>0.
\]
At $\tau=0$ the two roots are exactly $\pm d/2$; strict
convexity gives $D(0)<0$, even when $d$ is very small.
Consequently $D$ has at most one zero in this time interval.
At that zero $J_{uu}$ and $D'$ are positive and nonzero, so
it is a simple discriminant zero with the stated real-to-complex
crossing. The two roots cannot return without another zero of
$D$. Applying NN-A023 on its own local time interval identifies
the collision time; whenever both intervals apply they track
the same two zeros by their common initial sections and local
uniqueness. No condition on $\delta L$ or $\sqrt T L$ was used.
\end{proof}

\begin{theorem}[NN-A025: Poisson process of early CUE collisions]
\label{thm:analytic-cue-collision-process}
For Haar CUE, let $(a_{N,j},b_{N,j})$ be the real double-zero
collisions of $Z_{N,-a}$, with $a_{N,j}>0$ and real position
$b_{N,j}$ taken modulo $2\pi$. On any fixed rescaled time
window, with probability tending to one, all multiple-zero
events are such simple double collisions; each creates one
nonreal conjugate pair, and these pairs remain nonreal on that
window. The marked point process
\[
 \sum_j\delta_{(v_{N,j},\,b_{N,j}\bmod2\pi)},
 \qquad v_{N,j}=\frac{8N^{8/3}a_{N,j}}{(72\pi)^{2/3}},
\]
converges on bounded time windows to a Poisson point process
with intensity
\begin{equation}\label{eq:analytic-collision-ppp-intensity}
                   \frac32\sqrt v\,dv\,\frac{db}{2\pi}.
\end{equation}
In particular the number of nonreal conjugate pairs at time
$a=(72\pi)^{2/3}v/(8N^{8/3})$ converges in distribution to
$\operatorname{Poisson}(v^{3/2})$. The pair-count processes
converge in finite-dimensional distributions to the corresponding
inhomogeneous Poisson counting process. If $V_k$ is the limiting
$k$th collision time, then
\[
 V_k^{3/2}\sim\operatorname{Gamma}(k,1),\qquad
 f_{V_k}(v)=\frac{3}{2(k-1)!}
                       v^{3k/2-1}e^{-v^{3/2}}\quad(v>0).
\]
The limiting collision positions are independent uniform circular
marks, independent of the times. This is a CUE theorem; it
asserts no transfer to arithmetic zero statistics.
\end{theorem}

\begin{proof}
Fix $A>0$ and put
\[
 s_N=N^{-4/3},\quad T_N=A s_N^2,\quad
 \ell_N=N^{-7/6},\quad B_N=N^{1/24},\quad h_N=B_Ns_N.
\]
The three-point divided-difference estimate in NN-A024 gives
\[
 \mathbb P(\text{some circular arc of length }\ell_N
                        \text{ contains three zeros})
       \le C N^9\ell_N^8=O(N^{-1/3}).
\]
Work on the complementary event. Call each adjacent gap
$\delta\le h_N$ a small-pair cluster and each remaining root
a singleton. Since $2h_N<\ell_N$ eventually, no root belongs
to two small pairs. Every other root lies at distance at least
$\ell_N-h_N\ge\ell_N/2$ from a small pair. Two singleton
roots are at least $h_N$ apart. These assertions apply also
across the circular cut, by using periodic real lifts.

We need uniform control of the other periodic zeros. The
no-triples condition gives at most two zeros in each interval
of length $\ell_N$ on the periodic real line. Shell summation
therefore gives, for a pair center with its two selected roots
removed,
\begin{equation}\label{eq:analytic-cue-uniform-pair-SL}
 S_{\rm pair}\le C\ell_N^{-2},\qquad
 |L_{\rm pair}|\le
                  C\ell_N^{-1}\log(2/\ell_N).
\end{equation}
For the inverse-square sum, bound the shells at distances
$k\ell_N/2$ by $C(k\ell_N)^{-2}$ and sum $k^{-2}$,
including all periodic copies. For the logarithmic drift,
the original trigonometric product gives exactly
\[
 L_{\rm pair}=\frac12\sum_{j\ {
m outside\ pair}}
                         \cot((b-\theta_j)/2).
\]
The selected two cotangents cancel at their midpoint, as do
the two linear-factor logarithmic derivatives. On circular
distances $d\in(0,\pi]$, $|\cot(d/2)|\le C/d$; summing the
finitely many shells gives the logarithm in
\eqref{eq:analytic-cue-uniform-pair-SL}. This avoids assigning
a nonconvergent unsigned first trace to all periodic lifts.

For a singleton at $b$, its other zeros are at least $h_N$
away. At most a fixed number can be closer than $\ell_N$;
the remaining shells give
\begin{equation}\label{eq:analytic-cue-uniform-single-SL}
 S_{\rm single}\le C(h_N^{-2}+\ell_N^{-2}),\qquad
 |L_{\rm single}|\le
               C(h_N^{-1}+\ell_N^{-1}\log(2/\ell_N)).
\end{equation}
Here $L_{\rm single}=\frac12\sum_{j\ne i}
\cot((b-\theta_j)/2)$; the regular term from its own sine
factor vanishes at $b$.

Uniformly over all pair clusters,
\[
 (\delta^2+T_N)S_{\rm pair}
       \le C(h_N^2+T_N)\ell_N^{-2}=O(N^{-1/4})=o(1).
\]
Uniformly over all singletons,
\[
 T_NS_{\rm single}
       \le C_A(B_N^{-2}+N^{-1/3})=o(1).
\]
Lemma~\ref{lem:analytic-uniform-clusters} therefore applies
simultaneously to every cluster on $[0,T_N]$. Their moving
discs remain pairwise disjoint. Indeed pair discs have radii
$2\sqrt{\delta^2+T_N}=O_A(h_N)=o(\ell_N)$ and their center
shifts are at most $C_AT_N\ell_N^{-1}\log(2/\ell_N)
=o(\ell_N)$. Singleton discs have radius $2\sqrt{T_N}=O_A(s_N)$;
their shifts, by \eqref{eq:analytic-cue-uniform-single-SL}, are
\[
 O_A(T_N/h_N+T_N\ell_N^{-1}\log(2/\ell_N))=o(s_N).
\]
They are separated from one another by at least $h_N=B_Ns_N$,
and from any pair by at least $\ell_N/2$ before these vanishing
relative displacements. Each disc has diameter tending to zero,
so disjointness also holds modulo $2\pi$.

The finite polynomial in $w=e^{iz}$ from NN-A024 has degree
$N$, with nonzero constant and leading coefficients at every
finite $a$. The disjoint discs contain a total of exactly $N$
zeros modulo $2\pi$, counted with multiplicity. They therefore
exhaust the entire zero divisor throughout $[0,T_N]$. This is
the step that remains valid after some pairs become nonreal;
no sign inequality for the later global root flow is required.
The singleton roots stay real and simple. Each pair has exactly
its unique collision, if that time is in the window, and cannot
return during the window. In particular there are no other
early collisions or undetected nonreal roots.

For every small pair the uniform collision expansion gives
\[
 a_c=\frac{\delta^2}{8}(1+O(N^{-1/4})).
\]
Only gaps with $\delta=O_A(s_N)$ can thus collide before $T_N$.
For those gaps the sharper bound $\delta^2S=O_A(N^{-1/3})$
is available. The collision mark differs from the left endpoint
of its initial gap by at most
\[
 \delta/2+2T_N|L_{\rm pair}|+O(\delta(\delta^2S)^{3/2})=o(1).
\]
Consequently on each bounded rescaled time window the collision
process is, up to a uniformly vanishing displacement, the image
of the initial extreme-gap process under
\[
                  (y,b)\longmapsto(y^2,b),\qquad
                  y=N^{4/3}\delta/(72\pi)^{1/3}.
\]
Ben Arous--Bourgade, Theorem~1.1, printed p.~2650, states
that the process in the unnormalized coordinate
$u=N^{4/3}\delta$ converges to the Poisson process of intensity
$u^2du/(24\pi)\,db/(2\pi)$. Thus the $y$ coordinate has
intensity $3y^2dy\,db/(2\pi)$, and squaring gives
\eqref{eq:analytic-collision-ppp-intensity}.
Its mass on $[0,v]\times\mathbb R/(2\pi\mathbb Z)$ is
$v^{3/2}$. The limit has no atoms at deterministic time or
spatial boundaries. Tightness of the number of extreme gaps
in bounded windows, supplied by that same point-process
convergence, justifies the vanishing-displacement comparison
and all asserted count limits. The formula for $V_k$ follows
by the deterministic time change $v\mapsto v^{3/2}$ to a
rate-one Poisson process. The product intensity yields the
independent uniform marks.
\end{proof}

\paragraph{Primary source and scope.}
The imported random-matrix result is exactly Ben Arous--Bourgade,
\emph{Extreme gaps between eigenvalues of random matrices},
Annals of Probability~41 (2013), 2648--2681,
Theorem~1.1 p.~2650 and Corollary~1.2 p.~2651,
\url{https://math.nyu.edu/~bourgade/papers/ExtremeGaps.pdf},
DOI~10.1214/11-AOP710. The source's full marked intensity,
including $db/(2\pi)$, was checked. The new argument supplies
the heat evolution, exhaustion after the first collision,
and absence of recollisions on bounded extreme-time windows.
A, M and ordinary fixed-scale pair correlation do not supply
an arithmetic version of the imported extreme-gap process.
There is no expectation convergence or error rate for the
limiting Poisson law claimed here.

\paragraph{Geometric and representation meaning.}
The actual finite free family
\[
 \operatorname{Spec}\mathbb C[q,q^{-1},w,w^{-1}]/
       \left(\sum_{k=0}^Nc_kq^{(2k-N)^2}w^k\right)
                           \longrightarrow\mathbb G_{m,q},
                  \qquad q=e^{-a/4},
\]
is the family already identified in NN-A024. On each bounded
extreme-time interval, with probability tending to one, its
nearby analytic zero algebra splits into the tracked singleton
and rank-two cluster factors. Each rank-two factor has at
most one simple trace-discriminant zero there, with a quadratic
local normalization; each singleton remains etale. SGA~1,
Expos\'e~I, section~4, Proposition~10 is the exact finite-locally-free
trace-discriminant criterion, and section~3, Proposition~1 gives
the vanishing-differentials criterion. The resulting simple
branch points obey the proved marked Poisson law. Analytic
monodromy can be taken on a common small tubular parameter
neighborhood of the real interval, preserving the disjoint
cluster factors. Around the separate early branch points it
exchanges disjoint pairs of sheets, so these local transpositions
commute.
This is a concrete random complex algebraic discriminant family,
not an arithmetic Frobenius equidistribution theorem, a pure
motive, or a positive integral basis.
